\section{Pilot and Diagnostic Study}

This section specifies a diagnostic test of the failure mode that \method{}
targets, on the specific memory systems and task streams we evaluate; it introduces
no second method and makes no claim about self-evolving memory in general. Every
plotted quantity is an unresolved pilot output, and the study is governed by a
prespecified falsification gate.

\paragraph{Systems and streams.}
We reproduce two credit designs on a unified codebase---a heuristic-credit bank in
the style of ReasoningBank~\tbd{CITATION_REASONINGBANK} and a Monte-Carlo-learned
bank in the style of MemRL~\tbd{CITATION_MEMRL}---using the same backbone and
retriever. These are \emph{standardized credit-variant wrappers} that isolate the
credit signal on a shared codebase, not source-faithful re-implementations; we
report them as ``in the style of'' variants and do not attach the exact method names
to their numbers unless the baseline-fidelity check (Section~6) confirms
source-faithful behavior. We run them on two replayable streams (a text embodied
environment and a multi-hop retrieval task), which are required because
interventions must re-run paired rollouts from an identical initial state (AS2).
Stream length and seed count are fixed as a planned configuration.

\paragraph{Two-stage interventional protocol.}
This is the first appearance of \rit{}. Thresholds of scientific meaning are fixed
\emph{before} the pilot: the equivalence half-width $\delta$ (the smallest
memory-level contribution we treat as practically meaningful) and the minimal
effect of interest are pre-registered as user-approved values, not read from pilot
noise. Stage~1 is a pilot with $K=5$ paired rollouts that estimates variance;
Stage~2 then computes and freezes the formal $K$, the per-memory audit sample size
$n^{\text{audit}}$, and the seed count so the confidence-interval half-width falls
below the pre-fixed $\delta$. The confirmatory statistics use data disjoint from the
Stage-1 pilot.

At snapshot points $t \in \{100, 200, 300, 400, 500\}$ we sample retrieval events
and, for each audited memory, perform a paired leave-one-out comparison---the full
retrieved set versus the same set with the memory removed---holding context order
fixed and inserting a length-matched neutral pad whose neutrality is validated
separately (\tbd{PILOT_NEUTRAL_PAD_VALIDATION}). The pad makes each label a
\emph{pad-replacement} contribution; it may be interpreted as the literal removal
effect only if that pad-equivalence validation passes a pre-registered equivalence
test, otherwise the estimand stays pad-replacement and literal removal is left as
unvalidated future work (no no-pad arm is added in this study). The event label is
$\phirit_i(q)$, the $K$-rollout success-rate difference. To keep audit estimates
unbiased, the audit pool is a random or stratified probability sample with known
inclusion probabilities (any priority sampling is confined to the training pool),
and estimates are restricted to memories with adequate common support (coverage
\tbd{PILOT_AUDIT_SUPPORT_COVERAGE}). \rit{} labels are partitioned into four
disjoint roles---\emph{fit} (trains \aca{}), \emph{calibration} (isotonic
calibration), \emph{development} (provisional threshold/hyperparameter selection),
and a \emph{sealed final-audit} pool that scores CCC/SR/gates and remains blind to
model, threshold, calibration, and stopping selection; the split fractions are
locked before running (\emph{USER\_APPROVAL\_REQUIRED}). Rollout counts and snapshot
points are a planned configuration; the amortized cost is
\tbd{RIT_ROLLOUT_EQUIV_PER_EVENT} per
event, and the totals \tbd{PILOT_TOTAL_ROLLOUT_BUDGET} and
\tbd{PILOT_TOTAL_TOKEN_COST} appear in an appendix cost table.

\paragraph{Statistics.}
We separate the true memory-level aggregate target
$\phiagg_i = \E[\phicf_{i,t}]$ from its finite \rit{} audit estimate
$\phiaggrit_i$, reported with a task-clustered paired confidence interval.
Concretely, $\phiaggrit_i$ is the design-based (inverse-inclusion-probability
weighted) mean of the per-event labels $\phirit_i(q)$ over the probability-sampled
audit events of $m_i$; it is an estimate, not ground truth. All three memory-level
signals ($\coutil_i$, $\phibar_i$, and $\phiaggrit_i$) are computed on this same
audit population and context distribution, so their comparison is not confounded by
differing endogenous supports. The \emph{credit--contribution
correlation} at the memory level, $\mathrm{CCC}(s, \phiaggrit)$ for a ranking signal
$s \in \{\coutil, \phibar\}$, is the Spearman correlation between the signal and the
audit estimate. The \emph{superstition rate} $\mathrm{SR@}k(s)$ is the fraction of
the top-$k\%$ memories by $s$ whose contribution is non-positive, judged from the CI
of $\phiaggrit_i$ against the pre-fixed equivalence band $[-\delta, \delta]$---not a
bare $\le 0$ cut. We report $\mathrm{SR@}20\%$ split into a \emph{harmful} subset
($\phiaggrit_i$ confidently below $-\delta$) and a \emph{practically null} subset
(CI within $[-\delta, \delta]$); redundant-but-useful memories are separated and
treated in the AS3 boundary study. Figure~\ref{fig:diagnostic-baseline} plots the
baseline memory-level $\mathrm{CCC}$ and the $\mathrm{SR@}20\%$ trajectory over
deployment.

% Ledger IDs:
% [[TBD:PILOT_CCC_REASONINGBANK_T500]]
% [[TBD:PILOT_SR20_T100]]
% [[TBD:PILOT_SR20_T500]]
\begin{figure}[t]
\centering
\figureplaceholder{1.55in}{%
\textbf{Fig. 2 placeholder: Baseline diagnostic}\\[3pt]
(a) $x$: system co-occurrence utility $\coutil$; $y$: memory-level audit estimate $\phiaggrit$ (finite \rit{} output).\\
(b) $x$: task index $t$ (100--500); $y$: SR@20\%.\\
Methods: ReasoningBank and MemRL.\\
The axis ranges and statistical definitions are locked to Fig.~\ref{fig:mechanism-recovery}.}
\caption{Planned diagnostic of credit-contribution correlation and superstition rate over deployment. All plotted values remain unresolved pilot outputs.}
\label{fig:diagnostic-baseline}
\end{figure}


\paragraph{Planned readings and controls.}
We test whether the baseline memory-level correlation is low
(\tbd{PILOT_CCC_REASONINGBANK_T500} for the heuristic bank and
\tbd{PILOT_CCC_MEMRL_T500} for the learned bank) and whether the superstition
subsets rise over deployment: harmful from \tbd{PILOT_SR20_HARMFUL_T100} to
\tbd{PILOT_SR20_HARMFUL_T500} and practically null from \tbd{PILOT_SR20_NULL_T100}
to \tbd{PILOT_SR20_NULL_T500}, with the combined $\mathrm{SR@}20\%$ moving from
\tbd{PILOT_SR20_T100} to \tbd{PILOT_SR20_T500} (each a point estimate with its CI),
as predicted if correlational credit does not self-correct. Two controls guard
against alternative explanations. A split-half re-test of $\phirit$ gives
reliability \tbd{PILOT_SPLITHALF_PHI_RELIABILITY}, testing whether a low correlation
is a measurement artifact. Removing the entire positive-contribution subset and
measuring the success drop \tbd{PILOT_UTILITY_SUBSET_REMOVAL_DROP} tests whether
memory has value at all, so that a defective credit assignment---rather than
useless memory---is the operative problem.

\paragraph{Falsification gate.}
If the upper bound of the memory-level credit--contribution correlation interval
falls above the pre-registered bound (the source names $\mathrm{CCC} > 0.5$ as the
falsifier), or if the superstition subsets do not rise with deployment, the
attribution hypothesis is not supported; the decision uses the confidence interval,
not a point estimate. We will then narrow the claim to heterogeneous streams rather
than force a positive reading. The diagnostic yields two requirements for the method:
approximate the \rit{} credit online, and expose the same credit signal to
downstream consumers. Both are realized by \method{}, whose architecture is shown in
Figure~\ref{fig:method-architecture}.
